\chapter{Guía de Instalación, Uso y Demostración}\label{cap:guia}
Este capítulo constituye el manual de usuario del proyecto. Describe los requisitos del entorno, la instalación de las herramientas, la compilación y ejecución del repositorio oficial, el uso del motor complementario, la regeneración de los resultados y de este documento, y el guion recomendado para la demostración en la Feria Facultativa. Los comandos se presentan para un sistema GNU/Linux; en Windows y macOS son idénticos una vez instalado Rust.

\section{Requisitos del entorno}
El repositorio oficial solo requiere el compilador de Rust y su gestor Cargo. Las herramientas de documentación son opcionales y solo se necesitan para regenerar las gráficas y este documento. La Tabla~\ref{tab:requisitos} resume el software utilizado.

\begin{table}[!htbp]
\caption{Software requerido y opcional}\label{tab:requisitos}
\small
\begin{tabularx}{\textwidth}{@{}l l >{\raggedright\arraybackslash}X@{}}
\toprule
Herramienta & Uso & Necesidad \\
\midrule
Rust y Cargo (edición 2024) & Compilar y ejecutar el repositorio & Obligatoria \\
Git & Clonar el repositorio & Recomendada \\
Python 3 con \texttt{matplotlib} & Generar gráficas desde los CSV & Opcional (documentación) \\
TeX Live con XeLaTeX y \texttt{latexmk} & Compilar el PDF del documento & Opcional (documentación) \\
Pandoc y \texttt{python-docx} & Generar la versión DOCX editable & Opcional (documentación) \\
draw.io (diagrams.net) & Editar y exportar los diagramas & Opcional (documentación) \\
\bottomrule
\end{tabularx}
\fuente{Todas las herramientas son gratuitas y de código abierto. No se requiere GPU; la primera compilación descarga el crate \texttt{rand} desde crates.io.}
\end{table}

\section{Instalación}
Rust se instala con el instalador oficial \texttt{rustup}, que descarga el compilador, Cargo y la biblioteca estándar. A continuación se clona el repositorio oficial y se compila en modo optimizado.

\begin{lstlisting}[language={},numbers=none,caption={Instalación de Rust y obtención del repositorio oficial},label=cod:instalacion]
curl --proto '=https' --tlsv1.2 -sSf https://sh.rustup.rs | sh
source "$HOME/.cargo/env"
rustc --version && cargo --version
git clone https://github.com/lanorrron/neural-network.git
cd neural-network
cargo build --release
\end{lstlisting}

El archivo \texttt{Cargo.toml} declara una sola dependencia, \texttt{rand 0.9}, que Cargo descarga junto con sus dependencias transitivas en la primera compilación; después puede trabajarse sin conexión. Ningún crate de aprendizaje automático ni de álgebra lineal interviene en el proyecto.

\section{Ejecución de los experimentos}
El binario recibe como argumento el experimento que se desea ejecutar. La Tabla~\ref{tab:cli} describe las órdenes disponibles, definidas en \texttt{src/main.rs} (Código~\ref{cod:main}).

\begin{table}[!htbp]
\caption{Órdenes de la interfaz de línea de comandos del repositorio oficial}\label{tab:cli}
\small
\begin{tabularx}{\textwidth}{@{}l >{\raggedright\arraybackslash}X@{}}
\toprule
Orden & Efecto \\
\midrule
\texttt{cargo run} & Entrena la red de vocales e informa las cinco predicciones \\
\texttt{cargo run -- lineal} & Regresión $y=2x$ con una neurona lineal \\
\texttt{cargo run -- booleano} & Clasificación de 0 como falso y 1 como verdadero \\
\texttt{cargo run -- vocal <letra>} & Entrena la red de vocales y evalúa la vocal indicada (A, E, I, O o U) \\
\texttt{cargo run -- <letra>} & Forma abreviada de la orden anterior \\
\bottomrule
\end{tabularx}
\fuente{Una letra no válida produce un mensaje de error. Con \texttt{--release} se usa el binario optimizado.}
\end{table}

\begin{lstlisting}[language={},numbers=none,caption={Ejemplos de ejecución del repositorio oficial},label=cod:ejecucion]
cargo run --release -- lineal       # y = 2x, se detiene al alcanzar MSE <= 1e-12
cargo run --release -- booleano     # 250 epocas con gradientes por epoca
cargo run --release -- vocal a      # entrena y reconoce la vocal A
\end{lstlisting}

Los resultados se imprimen en la terminal; el programa no escribe archivos. Para documentar este trabajo, el script \path{scripts/auditar_repo_oficial.py} del repositorio de documentación compila la instantánea del commit \texttt{\OfCommit} guardada en \path{repo_oficial/neural-network/}, ejecuta las cuatro órdenes, guarda su salida en \path{resultados/oficial/*.txt} y la convierte en los CSV descritos en la Tabla~\ref{tab:csv}.

\begin{table}[!htbp]
\caption{Archivos de resultados de la auditoría del repositorio oficial}\label{tab:csv}
\small
\begin{tabularx}{\textwidth}{@{}l >{\raggedright\arraybackslash}X@{}}
\toprule
Archivo & Contenido \\
\midrule
\texttt{lineal.csv}, \texttt{lineal\_predicciones.csv} & Época, peso, bias y MSE; predicciones finales de $y=2x$ \\
\texttt{booleano.csv}, \texttt{booleano\_predicciones.csv} & Gradientes y parámetros por época; probabilidades finales \\
\texttt{vocales.csv}, \texttt{vocales\_predicciones.csv} & Probabilidad de cada vocal y cambio máximo por informe; predicciones finales \\
\texttt{vocal\_a.csv} & Probabilidades asignadas al evaluar la vocal A \\
\texttt{cargo\_test.txt}, \texttt{cargo\_clippy.txt} & Salida completa de las pruebas y del análisis estático \\
\texttt{resumen.csv} & Métricas finales, tiempos, advertencias y tamaño del binario \\
\bottomrule
\end{tabularx}
\fuente{Directorio \path{resultados/oficial/}. Auditoría completa en \path{docs/auditoria_repo_oficial.md}.}
\end{table}

\section{Verificación y motor complementario}
En el repositorio oficial, \texttt{cargo test} compila el proyecto y ejecuta \OfPruebas{} pruebas, y \texttt{cargo clippy} informa \OfClippyErr{} error (\texttt{modulo\_one}, en \texttt{vowels.rs}, porque \texttt{REPORT\_EVERY} vale 1) y \OfClippyWarn{} advertencias, en su mayoría código definido pero aún no utilizado. La verificación automatizada de las ecuaciones se realiza con el motor complementario, que se obtiene y ejecuta como sigue.

\begin{lstlisting}[language={},numbers=none,caption={Pruebas y experimentos del motor complementario},label=cod:pruebas]
git clone https://github.com/Noxy-bit/Redes-Neuronales-Rust.git
cd Redes-Neuronales-Rust
cargo test                          # 52 pruebas, incluido gradient checking
cargo run --release -- --all        # XOR, vocales con ruido y benchmarks -> resultados/*.csv
\end{lstlisting}

\section{Regeneración de gráficas, diagramas y documento}
Las gráficas, los diagramas y este documento se regeneran a partir de los resultados con los siguientes comandos. El script de compilación extrae los fragmentos de código, genera las macros de resultados desde los CSV, compila el PDF con XeLaTeX y produce la versión DOCX.

\begin{lstlisting}[language={},numbers=none,caption={Regeneración de la documentación},label=cod:regenerar]
python3 scripts/auditar_repo_oficial.py     # ejecuta el repositorio oficial -> resultados/oficial/
python3 scripts/graficas_repo_oficial.py    # 3 graficas oficiales (300 DPI)
python3 scripts/generar_graficas.py         # graficas del motor complementario
python3 diagramas/generar_diagramas.py      # diagramas .drawio -> PNG (requiere draw.io)
bash tesis/compilar_final.sh                # PDF y DOCX en tesis/build/
\end{lstlisting}

\section{Guion de demostración en la feria}
La demostración en el stand está pensada para durar entre cinco y siete minutos y puede adaptarse al tiempo disponible. Se recomienda tener el repositorio oficial compilado previamente y una terminal abierta en su raíz; las pausas del programa permiten seguir el entrenamiento en la pantalla.

\begin{enumerate}[leftmargin=2em,itemsep=4pt]
\item \textbf{Motivación (1 minuto).} Presentar el problema de la caja negra con el árbol de problemas (Figura~\ref{fig:arbol}) y la pregunta de investigación.
\item \textbf{Arquitectura (1 minuto).} Mostrar la Figura~\ref{fig:arquitectura}: vectores, neurona, error, gradiente y descenso de gradiente escritos a mano.
\item \textbf{Regresión en vivo (1 minuto).} Ejecutar \texttt{cargo run --release -- lineal} y observar cómo el peso se acerca a 2 y el MSE desciende hasta $10^{-12}$.
\item \textbf{Vocal elegida por el público (2 minutos).} Pedir a un visitante una vocal, ejecutar \texttt{cargo run --release -- vocal <letra>} y comentar cómo crece la probabilidad de la clase correcta; mostrar los patrones de la Figura~\ref{fig:of-voc}.
\item \textbf{Más allá de una capa (1 minuto).} Con el motor complementario, mostrar la frontera de XOR (Figura~\ref{fig:frontera}) y las pruebas de \emph{gradient checking}.
\item \textbf{Cierre (1 minuto).} Relacionar el proyecto con los contenidos de Programación 1 (Tabla~\ref{tab:inf120}) e invitar a escanear el código QR del repositorio oficial.
\end{enumerate}

\begin{nota}[Preguntas frecuentes del público]
\textbf{¿Por qué no usar Python y una biblioteca?} Porque el objetivo es comprender y verificar cada operación, no solo obtener un resultado. \textbf{¿La red aprende o memoriza?} La demostración de vocales usa un único patrón por letra; el motor complementario evalúa vocales con ruido no vistas en el entrenamiento y alcanza \EtresAccPru{} de exactitud. \textbf{¿Se necesita una computadora potente?} No: el binario oficial ocupa \OfBinKiB~KiB y su cómputo dura milisegundos; la mayor parte del tiempo en pantalla corresponde a las pausas de la demostración.
\end{nota}

\section{Solución de problemas frecuentes}
Si \texttt{cargo} no se reconoce como comando después de la instalación, debe abrirse una nueva terminal o ejecutarse \texttt{source "\$HOME/.cargo/env"}. Si la compilación indica que la edición 2024 no es compatible, debe actualizarse Rust con \texttt{rustup update stable}. Si la primera compilación falla al descargar \texttt{rand}, debe verificarse la conexión a Internet. Las métricas de los experimentos oficiales deben coincidir exactamente en cualquier equipo, porque no intervienen números aleatorios; los tiempos, en cambio, dependen del procesador. Si la generación de gráficas falla, debe instalarse \texttt{matplotlib} con \texttt{pip install matplotlib}.
